\section{条约法}
\begin{enumerate}
  \item 缔结。程序、保留。
        \begin{enumerate}
          \item 成立。
                \begin{enumerate}
                  \item 主体合法。国际法主体之间，\zd{不要缩小到国家}。
                        \begin{enumerate}
                          \item 全权证书\zd{总理 \hz 外交部长}签署，政府部门名义部门首长签。
                          \item 国家元首、政府首脑、外交部长、使馆馆长和派驻国际组织 \hz 会议的代表无须全权证书，\zd{正职}。
                        \end{enumerate}
                  \item 意思表示真实、一致。形式无要求，\zd{书面口头}均可。
                  \item 内容不违反国际强行法。
                  \item 缔约权不来自于公约，来自\zd{内部规则}。
                \end{enumerate}

                \begin{kk}{缔结程序}
                  \begin{enumerate}
                    \item 缔结新条约。
                          \begin{enumerate}
                            \item \zd{条约明确规定}其生效方式。如签署（书）、核准（书） \hz 批准（书）。
                            \item 没有规定 \yjt 重要类批准，非重要类核准。
                          \end{enumerate}
                    \item 进入老条约。\zd{开放性}条约才可进入。
                          \begin{enumerate}
                            \item 条约规定。如加入（书）、接受（书）。
                            \item 没有规定 \yjt 高规格加入，低规格接受。
                          \end{enumerate}
                    \item 重要类人常定国家主席签，其他国务院定外交部长签。\zd{核准书还可由总理签}。
                    \item \zd{谁定谁公布}。人常定人常公布，国务院定国务院公布。
                  \end{enumerate}
                  \tcblower
                  \zd{和好领界要批准，引渡协助法不同}。友好合作/和平、领土、司法协助/引渡、同我国法律有不同规定、各方议定须经批准的。
                \end{kk}
          \item \zd{外交部}在联合国\zd{秘书处}登记、公布，登记不是生效/产生约束力的方式，未经登记不得援引，
          \item 保留。是缔约国的权利，是\zd{单方声明}。
                \begin{enumerate}
                  \item 时间限制。\zd{尚未对保留国生效}才允许保留。
                  \item 接受和效力。\zd{规定禁止保留，则不得提出保留}。
                        \begin{enumerate}
                          \item 明确允许保留。
                          \item 缔约国数目有限且保留涉及条约宗旨需全体接受。
                          \item 国际组织章程，该组织有权机构接受。
                        \end{enumerate}
                        \begin{kk}{自行决定}
                          \begin{enumerate}
                            \item 同意。被保留的内容所取代。\al{不承认双重国籍}。
                            \item （小）反对。保留的条款\zd{视为不存在}。大反对。\al{丙反对整个条约在甲、丙之间生效}。
                          \end{enumerate}
                        \end{kk}
                \end{enumerate}
        \end{enumerate}
  \item 适用。
        \begin{enumerate}
          \item 对非缔约国创设义务/权利/取消 \yjt 第三国\zd{书面明示}接受/不反对既有效/原则上同意。
          \item 冲突。有约从约 \yjt 完全相同\zd{后约取代前约} \yjt 不完全相同 \yjt 个案处理。
          \item 解释。只能以\zd{作准文本}解释 \yjt \zd{善意解释}（违约有效性、中立、不占便宜）。
          \item 多边条约的修正。均接受用新的；不接受之间用老的；新加入的可以选，不选视为用新的。
          \item 违约的处理。Todo。
        \end{enumerate}
\end{enumerate}
